\begin{table}[!htbp]
\caption{Convergence characteristics of MS-AgentNet.}
\label{tab:4-3}
\raggedright
\zihao{-5}
\setlength{\tabcolsep}{2.4pt}
\renewcommand{\arraystretch}{1.04}
\begin{ThreePartTable}
\begin{adjustbox}{max width=\linewidth}
\begin{tabularx}{\linewidth}{>{\raggedright\arraybackslash}p{0.14\linewidth}>{\centering\arraybackslash}p{0.17\linewidth}>{\centering\arraybackslash}p{0.17\linewidth}>{\centering\arraybackslash}p{0.20\linewidth}>{\centering\arraybackslash}p{0.20\linewidth}}
\toprule
数据集 & 10\%阈值轮次中位数 & 阈值轮次范围 & 后期平均损失中位数 & 后期损失标准差中位数 \\
\midrule
Oxford & 26 & 18–35 & $8.14\times10^{-4}$ & $2.38\times10^{-4}$ \\
MIT & 7 & 6–8 & $5.84\times10^{-5}$ & $4.84\times10^{-6}$ \\
\bottomrule
\end{tabularx}
\end{adjustbox}
\begin{tablenotes}[flushleft]
\footnotesize
\item[] \parbox[t]{0.98\linewidth}{Note: Statistics are calculated from runs using five different random seeds. The threshold-epoch range is reported as the minimum--maximum.}
\end{tablenotes}
\end{ThreePartTable}
\end{table}
